\subsection{\texorpdfstring{DISCRETE SUBGROUPS OF \(\mathbf{R}^n\)}{DISCRETE SUBGROUPS OF R TO THE N}}

We have seen in Chapter V (§ 1, no. 1, Proposition 1) that the only closed subgroups of R, other than R itself, are the discrete subgroups, generated by a single element. We shall begin by considering the discrete subgroups of \(\mathbf{R}^n\) .

First of all, the subgroup of \(\mathbf{R}^n\) generated by \(p\) vectors ( \(p \leqslant n\) ) of the canonical basis (Chapter VI, § 1, no. 3) of \(\mathbf{R}^n\) is a discrete group isomorphic to the product \(\mathbf{Z}^p\) of \(p\) groups which are equal to \(\mathbf{Z}\) . More generally, consider the subgroup \(\mathbf{G}\) generated by \(p\) points \(a_i\) ( \(i \leqslant i \leqslant p\) ) which form a free system. There is a bijective linear mapping of \(\mathbf{R}^n\) onto itself which maps \(a_i\) to \(e_i\) ( \(i \leqslant i \leqslant p\) ); since such a mapping is an automorphism of the topological group \(\mathbf{R}^n\) , \(\mathbf{G}\) is isomorphic as a topological group to the subgroup generated by the \(e_i\) ( \(i \leqslant i \leqslant p\) ) and is therefore a discrete subgroup of rank \(p\) isomorphic to \(\mathbf{Z}^p\) .

The structure of the group \(Z^{p}\) , and hence that of G, has been studied in algebra. We recall the main results of this study. The bases of G with respect to the ring Z are systems of p points

\[
\boldsymbol {b} _ {i} = \sum_ {j = 1} ^ {p} r _ {i j} \boldsymbol {a} _ {j},
\]

where the \(r_{ij}\) are integers such that the determinant \(\det(r_{ij})\) is equal to +1 or -1. Every subgroup H of G is discrete and of rank \(q \leqslant p;\) furthermore, if \(H\) is a given subgroup of rank \(q\) , there exists a free system of \(p\) points \(b_i\) ( \(i \leqslant i \leqslant p\) ) which generates \(G\) , and a system of \(q\) points \(c_i\) ( \(i \leqslant i \leqslant q\) ) which generates \(H\) , such that we have \(c_i = e_i b_i\) for \(i \leqslant i \leqslant q\) , where the \(e_i\) are integers (the invariant factors of \(H\) with respect to \(G\) ) such that

\[
e _ {i + 1} \equiv o (\mathrm{mod} e _ {i}) \quad \text { for } \quad i \leqslant i \leqslant q - i.
\]

The quotient group G/H is a discrete group isomorphic to \(Z^{p-q} \times F\) , where F is a finite abelian group, the direct product of q cyclic subgroups of respective orders \(e_{1}, \ldots, e_{q}\) .

We shall show now that the discrete subgroups of \(\mathbf{R}^n\) which we have just been considering are the only ones that exist.

PROPOSITION 1. Let G be a discrete subgroup of \(\mathbf{R}^n\) of rank \(p\) , let \((\boldsymbol{a}_i)_{1 \leqslant i \leqslant p}\) be a free system of \(p\) points of G, and let P be the closed parallelotope with centre o and basis vectors \(\boldsymbol{a}_i\) (Chapter VI, § 1, no. 3). Then the set G ∩ P is finite and generates G, and every point of G is a linear combination of the \(\boldsymbol{a}_i\) with rational coefficients.

G ∩ P is compact and discrete, hence finite. Let x be any point of G; it is equal to a linear combination \(\sum_{i=1}^{p} t_{i}a_{i}\) of the \(a_{i}\) with real coefficients. For each integer m \textgreater{} 0, consider the point

\[
z _ {m} = m x - \sum_ {i = 1} ^ {p} [ m t _ {i} ] a _ {i} = \sum_ {i = 1} ^ {p} (m t _ {i} - [ m t _ {i} ]) a _ {i} (*);
\]

it belongs to G, and since \(o \leqslant mt_{i} - [mt_{i}] < 1\) , it lies in P. Hence, first, \(x = z_{1} + \sum_{i=1}^{p} [t_{i}] a_{i}\) , so that G is generated by \(G \cap P\) ; and secondly, since \(G \cap P\) is finite, there exist two distinct integers h, k such that \(z_{h} = z_{k}\) , so that \((h - k)t_{i} = [ht_{i}] - [kt_{i}]\) less \((1 \leqslant i \leqslant p)\) , and therefore the \(t_{i}\) are rational numbers.

COROLLARY. Let \((\pmb{a}_i)_{1\leqslant i\leqslant p}\) be a free system of \(p\) points of \(\mathbf{R}^n\) and let

\[
\boldsymbol {b} = \sum_ {i = 1} ^ {p} t _ {i} \boldsymbol {a} _ {i}
\]

be a linear combination of the \(a_{i}\) , with real coefficients. Then the subgroup G of \(R^{n}\) generated by the \(p + i\) points \(a_{1}, a_{2}, \ldots, a_{p}\) and b is discrete if and only if the numbers \(t_{i}\) are rational.

Proposition 1 shows that the condition is necessary. Also it is sufficient, for if it is satisfied we can write \(t_i = m_i / d\) , where \(d\) and the \(m_i\) are integers ( \(i \leqslant i \leqslant p\) ); \(b\) is therefore a linear combination, with integer coefficients, of the \(p\) points ( \(\mathrm{I} / d\) ) \(a_i\) ; hence \(G\) is a subgroup of the discrete group generated by these \(p\) points, and so is itself discrete.

The result of Proposition 1 can be expressed as follows: if q points \(x_{i}\) ( \(i \leqslant i \leqslant q\) ) of a discrete subgroup G of \(R^{n}\) form a system which is dependent with respect to R, then they form a system which is dependent with respect to Q. It follows immediately that the rational rank of a discrete subgroup of \(R^{n}\) is equal to its rank.

The Corollary to Proposition 1, applied to the case where the \(a_{i}\) are the \(n\) vectors \(\pmb{e}_{i}\) of the canonical basis of \(\mathbf{R}^{n}\) , gives us the following proposition:

PROPOSITION 2 (Kronecker). Let \(\theta_{1}, \theta_{2}, \ldots, \theta_{n}\) be \(n\) real numbers. In order that, for each \(\varepsilon > 0\) , there exist an integer \(q\) and \(n\) integers

\[
p _ {i} \quad (\mathrm{I} \leqslant i \leqslant n)
\]

such that

\[
\left| q \theta_ {i} - p _ {i} \right| \leqslant \varepsilon \quad (\mathrm{I} \leqslant i \leqslant n),
\]

where the left-hand side of at least one of these inequalities does not vanish, it is necessary and sufficient that at least one of the \(\theta_{i}\) be irrational.

THEOREM 1. Every discrete subgroup G of \(\mathbf{R}^n\) , of rank equal to \(p\) , is generated by a free system of \(p\) points.

By virtue of the properties of groups isomorphic to \(Z^{p}\) recalled earlier, it is enough to show that G is a subgroup of a discrete group generated by a free system of p points. Now since G is of rank p there exists a free system of p points \(a_{i}\) ( \(1 \leqslant i \leqslant p\) ) of G such that every \(x \in G\)

is equal to a linear combination \(\sum_{i=1}^{p} t_i a_i\) of the \(a_i\) with real coefficients;

since G is discrete, Proposition 1 shows that the \(t_i\) are rational. Furthermore, Proposition 1 shows that G is generated by a finite number of points; since these points are linear combinations of the \(\boldsymbol{a}_i\) with rational coefficients, there exists an integer \(d\) such that they are linear combinations with integer coefficients of the \(p\) points \((\mathrm{I}/d)\boldsymbol{a}_i = \boldsymbol{a}_i'\) . It follows that G is a subgroup of the group generated by the \(\boldsymbol{a}_i'\) .

Theorem 1 can also be proved without recourse to the theory of invariant factors (cf.~Exercise 1).

Discrete subgroups of \(R^{n}\) which are of rank n are also called lattices in \(R^{n}\) .

